\documentclass[a4paper]{article}
\usepackage{fontspec}
\usepackage{polyglossia}
\setdefaultlanguage{vietnamese}
\usepackage{mathptmx}
\setmainfont{Times New Roman}
\usepackage{geometry,graphicx,fancyhdr,lastpage,hyperref,caption,subcaption}
\usepackage{booktabs,longtable,array,tabularx,multirow,diagbox}
\usepackage{amsmath,amssymb,amsthm}
\usepackage{xcolor,listings,tikz,mdframed,algorithm2e}
\geometry{a4paper,total={160mm,240mm},left=30mm,top=22mm,right=20mm,bottom=20mm}
\newtheorem{theorem}{Định lý}
\begin{document}
\section{Kiểm tra ngoại tuyến}
Tiếng Việt: Nguyễn Thị Ánh, Trường Đại học Bách khoa.
Tệp nguồn: \texttt{BTL\_HK232.pdf}.

\begin{theorem}
Với mọi $a,b \in \mathbb{R}$, $(a+b)^2=a^2+2ab+b^2$.
\end{theorem}

\subsection{Công thức ở cỡ chữ tiêu đề $h=10$}
Giá trị $P_{2000}=6$ phải biên dịch được hoàn toàn ngoại tuyến.

\begin{algorithm}[H]
\KwIn{Một tài liệu Markdown}
\KwOut{Một tệp PDF}
Biên dịch tài liệu\;
\end{algorithm}

\begin{lstlisting}
print("offline")
\end{lstlisting}

\begin{tikzpicture}
  \draw (0,0) rectangle (2,1);
\end{tikzpicture}

\begin{tabular}{ll}
\toprule
Hạng mục & Trạng thái \\
\midrule
Bundle & Ngoại tuyến \\
\bottomrule
\end{tabular}
\end{document}
